\SetPractica{Dureza total}{Determinación de dureza total de agua}

\section{Introducción teórica}

La dureza se entiende como la capacidad de un agua para precipitar al jabón y esto está basado en la presencia de sales de los iones calcio y magnesio. La dureza es la responsable de la formación de incrustaciones en recipientes y tuberías lo que genera fallas y pérdidas de eficiencia en diferentes procesos industriales como las unidades de transferencia de calor. El término dureza se aplicó en principio por representar al agua en la que era difícil (duro) de lavar y se refiere al consumo de jabón para lavado, en la mayoría de las aguas alcalinas esta necesidad de consumo de jabón está directamente relacionada con el contenido de calcio y magnesio.

El método se basa en la formación de complejos por la sal disódica del ácido etilendiaminotetraacético con los iones calcio y magnesio. El método consiste en una valoración empleando un indicador visual de punto final, el negro de eriocromo T, que es de color rojo en la presencia de calcio y magnesio y vira a azul cuando estos se encuentran acomplejados o ausentes. El complejo del EDTA con el calcio y el magnesio es más fuerte que el que estos iones forman con el negro de eriocromo T, de manera que la competencia por los iones se desplaza hacia la formación de los complejos con EDTA desapareciendo el color rojo de la disolución y tornándose azul.

\section{Objetivos}

\begin{itemize}
    \item 
\end{itemize}

\section{Materiales, equipos y reactivos}

\begin{table}[H]
\centering{\small\setstretch{1.25}
\begin{tabular}{|m{0.3\linewidth}|m{0.3\linewidth}|m{0.3\linewidth}|} \hline
    
    \thead{Equipos} & \thead{Materiales} & \thead{Reactivos} \\ \hline
    
    {Agitador magnético \par}
    {Imán de agitación \par}
    {Potenciómetro}
    
    &
    
    {Pipeta graduada \par}
    {Gotero \par}
    {Probeta \par}
    {Vasos de precipitado \par}
    {Matraz Erlenmeyer \par}
    {Soporte universal \par}
    {Pinza para bureta \par}
    {Bureta \par}
    {Colador}
    
    & 
    
    {EDTA 0.02M \par}
    {Buffer para dureza \par}
    {Negro de ericromo T}
    
    \\ \hline
\end{tabular}}\end{table}

\section{Procedimientos}
\begin{enumerate}
    \item Llene la bureta con el EDTA 0.01 M.
    \item Tome 50mL de agua de la llave del baño
    \item Agregue 2mL de buffer para dureza en un erlenmeyer de 125mL
    \item Agregue una pizca de negro de eriocromo T como indicador
    \item Titule hasta la aparición del cambio de color rosado a azul.
    \item Anote el volumen de titulante gastado.
    \item Realice este proceso al menos 2 veces
    \item Calcular el porcentaje de acidez de acuerdo con la siguiente fórmula:
    \begin{equation*}
        \frac{mg}{L} = \frac{ M_{EDTA} \cdot V(mL)_{EDTA \cdot \frac{a}{b} \cdot PF(CaCoO_{3}) } }{ V_{\text{muestra en L}} }
    \end{equation*}
    
    \item Repetir con las otras muestras 
\end{enumerate}